\section{Sources and Rights}
\label{sec:appendix-sources}

\subsection{The three-layer registry}

Rights are recorded at three granularities that are allowed to disagree: the host
(22 records), the individual file (53 artefacts, each with its own licence field),
and the text layer (15 records). The governing statement is that rights attach to
each source file, translation, transcription, image or recording, and that website
accessibility does not grant reuse.

Thirteen numbered policy rules govern reconciliation. The load-bearing ones: on
conflict the more restrictive statement controls; a permissive status is never
inferred, only evidenced; absence of a licence field is not a grant; chain of title
is followed to its root rather than to the nearest licence; a digitiser's stamp is
never a grant and a printed notice controls; suitability for one Veda does not
transfer to another; and audio carries its own rights regardless of the text's.

Status distribution at file level: \code{UNKNOWN} 15, \code{PUBLIC\_DOMAIN} 10,
\code{APACHE\_2\_0} 10, \code{PERMISSION\_REQUIRED} 6, \code{CC\_BY\_SA} 6,
\code{REFERENCE\_ONLY} 3, \code{CC\_BY\_NC\_SA} 2, \code{CC\_BY} 1.

\subsection{Text versions actually held}

\begin{center}\small
\begin{tabularx}{\linewidth}{@{}l>{\raggedright\arraybackslash}Xl@{}}
\toprule
Text version & Role & Rights \\
\midrule
\code{GRETIL.RV.AUFRECHT} & primary Sanskrit, \d{R}gveda & \code{CC\_BY\_NC\_SA} \\
\code{VEDAWEB.AUFRECHT} & parallel reading & \code{CC\_BY\_NC\_SA} \\
\code{VEDAWEB.EICHLER} & parallel, accented Devanagari & \code{CC\_BY\_NC\_SA} \\
\code{VEDAWEB.VNH} & metrically restored; \emph{never} primary & \code{CC\_BY\_NC\_SA} \\
\code{VEDAWEB.LUBOTSKY} & comparison only & \code{CC\_BY\_NC\_SA} \\
\code{VEDAWEB.PADAPATHA} & declared, \textbf{never ingested} & \code{CC\_BY\_NC\_SA} \\
\code{VEDAWEB.ZURICH} & linguistic annotation & \code{CC\_BY} \\
\code{WIKISOURCE\_SA.SV.KAU.ARCIKA\_MULA} & primary Sanskrit, S\=amaveda & \code{CC\_BY\_SA} \\
\code{GRETIL.SV.KAUTHUMA} & primary; \textbf{rejected witness} & \code{PERMISSION\_REQUIRED} \\
\code{WIKISOURCE\_SA.YV.VSM.ACCENTED} & extracted from container & \code{CC\_BY\_SA} \\
\code{WIKISOURCE\_SA.YV.VSM.UNACCENTED} & parallel reading & \code{CC\_BY\_SA} \\
\code{GRETIL.AVS.SAUNAKA.ACCENTED} & primary Sanskrit, Atharvaveda & \code{REFERENCE\_ONLY} \\
\code{GRETIL.AVS.SAUNAKA.UNACCENTED} & parallel reading & \code{REFERENCE\_ONLY} \\
\code{VEDAGRAPH.AVS.ROTH\_WHITNEY\_1856.TRANSCRIPTION} & the \textbf{cancelled}
  independent transcription & \code{CC0} \\
\code{VEDAGRAPH.AVS.SEARCH\_NORMALIZED} & derived search surface & \code{REFERENCE\_ONLY} \\
\bottomrule
\end{tabularx}
\end{center}

Three rows in that table are not data the project uses, and they are listed because
their presence is the point: a declared-but-never-ingested Padap\=a\d{t}ha, a
rejected S\=amavedic witness retained as the record of a defect family, and a
cancelled transcription programme.

\subsection{Printed editions behind the digital artefacts}

The \d{R}gvedic primary text descends from Aufrecht's edition
\citep{aufrecht1877rigveda} through a data-entry, conversion and \mbox{TEI} chain
that the artefact registry records step by step, with a pinned checksum and a
publication date for the \mbox{TEI} encoding \citep{gretil}. That file's header does
not state a recension, so the \'S\=akala identification is recorded as a reviewed
structural inference rather than as a source statement.

Traditional \d{R}gvedic metadata comes from a digital Anukrama\d{n}\=\i\
\citep{akavarapu2023anukramani}, commit-pinned, with a four-layer separation recorded
explicitly: the traditional source is cited and not ingested; the digital
transcription is ingested; modern editorial normalisation is recorded and never
undone; and machine-classification labels in the same repository are never ingested
at all.

English renderings are Griffith for the \d{R}gveda, Yajurveda and Atharvaveda
\citep{griffith1896rigveda,griffith1899yajurveda,griffith1895atharvaveda} and Whitney
with Lanman for the Atharvaveda \citep{whitney1905atharva}. Both are public domain;
the project depends on the public-domain status of the \emph{translation text} and
deliberately not on the transcription layer carrying it, which is separately
licensed.

\subsection{The combined-release constraint}

The \d{R}gvedic primary text is CC~BY-NC-SA, the S\=amavedic and Yajurvedic texts are
CC~BY-SA, and the Atharvavedic text is reference-only. The registry's own analysis is
that under CC~BY-SA's share-alike clause a combined four-Veda release is not
licensable as a single adapted work from these sources. This is why the graph is
buildable and queryable but not redistributable as one artefact
(\cref{sec:data-availability}).

\subsection{Snapshot provenance}

An independent audit over the raw store verified 1{,}347 metadata records against
their payloads: 1{,}347 hash-verified, 0 mismatches, 0 missing payloads, 0 orphans,
0 incomplete provenance records and 0 partial-download residue, across roughly
249\,MB. Eleven sidecar fields are required per snapshot, including retrieval URL,
retrieval timestamp, HTTP status, request headers and parser-independent metadata.
A repository URL without a commit pin raises.

The audit states its own limits, and we adopt them: hash verification proves a
snapshot is unaltered since retrieval and says nothing about whether retrieval was
permitted, nor about whether the digest was pinned against the right artefact.
Integrity and scope are independent properties.
